% concatenated from PaperHp.tex
\documentclass[a4paper,12pt,reqno]{amsart}
\usepackage{eurosym}
\usepackage[T1]{fontenc}
\usepackage{amsthm}
\usepackage{amsmath}
\usepackage{amssymb}
\usepackage{mathrsfs}
\usepackage{latexsym}
\usepackage{exscale}
\usepackage{geometry}
\usepackage{graphicx}
\usepackage[utf8]{inputenc}
\usepackage{color}
\usepackage[colorlinks=true, pdfstartview=FitV, linkcolor=blue, citecolor=red, urlcolor=blue]{hyperref}

\setcounter{MaxMatrixCols}{10}
%TCIDATA{OutputFilter=LATEX.DLL}
%TCIDATA{Version=5.50.0.2953}
%TCIDATA{<META NAME="SaveForMode" CONTENT="1">}
%TCIDATA{BibliographyScheme=Manual}
%TCIDATA{LastRevised=Tuesday, July 19, 2022 10:08:19}
%TCIDATA{<META NAME="GraphicsSave" CONTENT="32">}
%TCIDATA{Language=American English}

\definecolor{red}{rgb}{1,0.1,0.1}
\definecolor{blue}{rgb}{0.1,0.1,1}
\definecolor{vb}{RGB}{160,32,240}
\newcommand{\red}{\color{red}}
\newcommand{\blue}{\color{blue}}
\newcommand{\vb}{\color{vb}}
\numberwithin{equation}{section}
\newtheorem{theorem}{Theorem}[section]
\newtheorem{acknowledgement}[theorem]{Acknowledgement}
\newtheorem{algorithm}[theorem]{Algorithm}
\newtheorem{axiom}[theorem]{Axiom}
\newtheorem{case}[theorem]{Case}
\newtheorem{claim}[theorem]{Claim}
\newtheorem{conclusion}[theorem]{Conclusion}
\newtheorem{condition}[theorem]{Condition}
\newtheorem{conjecture}[theorem]{Conjecture}
\newtheorem{corollary}[theorem]{Corollary}
\newtheorem{criterion}[theorem]{Criterion}
\newtheorem{definition}[theorem]{Definition}
\newtheorem{example}[theorem]{Example}
\newtheorem{exercise}[theorem]{Exercise}
\newtheorem{lemma}[theorem]{Lemma}
\newtheorem{notation}[theorem]{Notation}
\newtheorem{problem}[theorem]{Problem}
\newtheorem{proposition}[theorem]{Proposition}
\newtheorem{remark}[theorem]{Remark}
\newtheorem{solution}[theorem]{Solution}
\newtheorem{summary}[theorem]{Summary}
\headheight=8pt
\topmargin=0pt
\textheight=659pt
\textwidth=436pt
\oddsidemargin=10pt
\evensidemargin=10pt
\calclayout
\allowdisplaybreaks
\newcommand{\NN}{\mathbb{N}}
\newcommand{\RR}{\mathbb{R}}
\newcommand{\GG}{\mathbb{G}}
\newcommand{\Ll}{\mathcal{L}}
\newcommand{\Xx}{\mathcal{X}}
\newcommand{\BB}{\mathcal{B}}
\newcommand{\Div}{\text{div}}
\newcommand{\Ff}{\mathcal{F}}
\DeclareMathOperator*{\esup}{ess\,sup}
\DeclareMathOperator*{\einf}{ess\,inf}


\title{Weighted $H^p\!- \! L^q$ boundedness  of  integral operators with rough kernels}

\author[A. L. Gallo]{Andrea L. Gallo}
\address{A. L. Gallo\\
	FaMAF \\
	Universidad Nacional de C\'ordoba \\
	CIEM (CONICET) \\
	5000 C\'ordoba, Argentina}
\email{andrea.gallo@unc.edu.ar}

\author[M. S. Riveros]{M. Silvina Riveros}
\address{M. S. Riveros\\
	FaMAF \\
	Universidad Nacional de C\'ordoba \\
	CIEM (CONICET) \\
	5000 C\'ordoba, Argentina}
\email{sriveros@unc.edu.ar}

\author[L. A. Vallejos]{Lucas A. Vallejos}
\address{L. A. Vallejos\\
	FaMAF \\
	Universidad Nacional de C\'ordoba \\
	CIEM (CONICET) \\
	5000 C\'ordoba, Argentina}
\email{lucas.vallejos@unc.edu.ar}


\thanks{The  authors are partially supported
	by CONICET and SECYT-UNC}

\subjclass[2000]{42B20, 42B25, 42B30}

\keywords{Fractional operators,  weighted   Hardy spaces, Muckenhoupt weights.}

\date{July 2026}

\begin{document}
	
	\maketitle
	
	\begin{abstract}
		In this paper, we study integral operators 
		\begin{equation*}
			T_\alpha f(x)=\int_{\mathbb{R}^{n}}K(x,y) f(y)dy,
		\end{equation*}
		with kernels $K(x,y)= k_1( x- A_1y)...k_m( x-A_my),$
		where $k_i(x)=\frac{\Omega_i(x)}{|x|^{n/q_i}}$ and $\Omega_i: \mathbb{R}^n\to \mathbb{R}$  are  homogeneous functions of degree zero, satisfying a size and a Dini condition, $A_{i}$ are certain invertible matrices, and  \   $\frac n{q_1}+\dots + \frac n{q_m}=n-\alpha,$  $0\leq \alpha <n.$
		We obtain the $H^{p}_{w^p}(\mathbb{R}^{n})-L^{q}_{w^q}(\mathbb{R}^{n})$  boundedness of these operators, for a class of Muckenhoupt weights $w$, satisfying the condition 
		\begin{equation*}
			w(A_ix)\leq cw(x),
		\end{equation*}
		\emph{ a.e.} $x\in\mathbb{R} ^n $, $1\leq i\leq m$.
	\end{abstract}
	
	
	\section{Introduction}
	
	The Hardy spaces on $\mathbb{R}^n$ were introduced in \cite{FS} by C. Fefferman and E. Stein, and have been extensively studied ever since. 
	A key feature of these spaces is that they provide an effective alternative to Lebesgue spaces in the range $0<p\leq 1$. 
	In particular, while Riesz transforms fail to be bounded on $L^p(\mathbb{R}^n)$ for $p\leq 1$, they do exhibit boundedness when considered on the corresponding Hardy spaces $H^p(\mathbb{R}^n)$.
	To investigate the boundedness of operators such as singular integrals or fractional-type operators on Hardy spaces $H^p(\mathbb{R}^n)$, one typically relies on their atomic or molecular characterizations. 
	This means that any distribution in $H^p$ can be expressed as a sum of atoms or molecules. 
	The atomic decomposition was first established by Coifman in one dimension (see \cite{C}) and later, in \cite{L}, extended to higher dimensions by Latter. 
	
	These descriptions suggest that the boundedness of linear operators on $H^p$ could, in principle, be reduced to analyzing their action on atoms or molecules. However, this approach has certain limitations. 
	Indeed, Bownik in \cite{B}, building on an example due to Meyer, constructed a linear functional defined on a dense subspace of $H^1(\mathbb{R}^n)$ that maps all $(1,\infty,0)-atoms$ into bounded scalars, but does not admit a bounded extension to the whole space $H^1(\mathbb{R}^n)$. This shows that, in general, it is not sufficient to verify that an operator sends atoms into bounded elements of a quasi-Banach space $X$ to guarantee that it extends to a bounded operator on $H^p$, for $0<p\leq 1$. 
	
	Nevertheless, such examples can be regarded as somewhat pathological. 
	For classical operators $T$, uniform boundedness on atoms implies boundedness from $H^p$ into $L^p$. 
	This follows from the fact that $T$ is bounded on $L^s$, for $1<s<\infty$, together with the possibility of choosing atomic decompositions that converge in the $L^s$ norm, see for example the papers of D.\@ Yang and Y.\@ Zhou \cite{YZ} and P.\@ Rocha \cite{Rocha-2}.
	
	The weighted Lebesgue spaces $L^p_{w}(\mathbb{R}^n)$ extend the classical Lebesgue spaces $L^p(\mathbb{R}^n)$ by replacing the Lebesgue measure $dx$ with the weighted measure $w(x)dx$, where $w$ is a non-negative measurable function. 
	In this setting, one can introduce the weighted Hardy spaces $H^p_{w}(\mathbb{R}^n)$ by adapting the definition of the classical Hardy spaces $H^p(\mathbb{R}^n)$ (see \cite{ST}). 
	It is well-known that harmonic analysis in this framework becomes particularly meaningful when the weight $w$ belongs to the Muckenhoupt class $A_{\infty}$.
	
	The atomic decomposition of $H^p_{w}(\mathbb{R}^n)$ was established in earlier contributions (see \cite{GC}, \cite{ST}). 
	Boundedness results for classical singular integral operators on $H^p_{w}(\mathbb{R}^n)$ were obtained under the assumption that $w \in A_1$. 
	Also, in \cite{Rocha-3}, the author extended these results for all $w \in A_{\infty}$. 
	Using these results, the author obtained the boundedness of certain singular integral operators on $H^p_w$ and from $H^p_w$	into $L^p_w$, for all weights $w \in A_{\infty}$ and $0<p \leq 1$. In addition, the author obtained the boundedness of the Riesz potential $I_{\alpha}$ from $H^p_w$ into $H^q_w$ where $0<p \leq 1$ and $\frac{1}{q}=\frac{1}{p} - \frac{\alpha}{n}$ and $w$ satisfying some appropriate conditions.
	
	
	Let $\Omega \in L^r(S^{n-1})$ be a homogeneous function of degree zero, where $r>1$ and $S^{n-1}$ denotes the unit sphere in $\mathbb{R}^n$, $n\geq 2$.
	For integral operators with rough kernels of the form	
	$$
	T_{\Omega, \alpha} f(x)=\int  \frac{\Omega(x-y)}{|x-y|^{n-\alpha}} f(y) \, dy,
	$$
	in \cite{KuW}, \cite{Du} and \cite{Wa}  the authors obtained weighted estimates for $T_{\Omega, 0}$ (in the principal  value sense) with $\Omega$ satisfying some additional conditions.  
	In \cite{DL}, the authors showed the corresponding weighted results for $\alpha>0$. 
	Also, in \cite {BLRi}, the authors obtained a Coifman type inequality for general fractional integral operators with kernels satisfying a H\"{o}rmander condition given by a certain Young function.
	
	In \cite{DL2}, Y.\@ Ding and S.\@ Lu applied the atomic decomposition and the molecular characterization of the real Hardy space to give the $H^p(\mathbb{R}^n)-L^q(\mathbb{R}^n)$ and  $H^p(\mathbb{R}^n)-H^q(\mathbb{R}^n)$ boundedness of $T_{\Omega, \alpha}$, for $0< p\leq 1$. In \cite{SW}, J.\@ O.\@ Str\"{o}ngberg and R.\@ L.\@ Wheeden gave the weighted $H^p-L^q$ and $H^p-H^q$ boundedness of the Riesz potencial $I_\alpha$.
	In \cite{DLeLi}, applying the atom-molecule theory developed  by J.\@ Garc\'{\i}a Cuerva, M.\@ Lee and C.\@ Lin (see \cite{GC} and \cite{LeLi}), Y.\@ Ding  M.\@ Lee and C.\@ Lin obtained the weighted  $H^p-L^q$ and $H^p-H^q$ boundedness of $T_{\Omega, \alpha}$.
	
	Let $0\leq \alpha<n$ and  $m\in \mathbb{N}$ with $m>1$. For  $1\leq i\leq m$,  let  $1<q_i<\infty$ such that 
	\begin{equation}\label{q_i}
		\frac n{q_1}+\dots+\frac n{q_m}=n-\alpha.
	\end{equation}
	Let $\Omega _i\in L^1(S^{n-1})$. If
	$x\neq 0$, we write $x'=x/|x|$. We extend this function to $\mathbb{R}^n\setminus\{0\}$ by $\Omega_i(x)=\Omega_i(x')$.
	%$\Omega_i$ is homogeneous of degree $0$.
	Let  \begin{equation}\label{ki} k_i(x)=\frac{\Omega_i(x)}{|x|^{n/q_i}}.\end{equation}
	In this paper, we  study the integral operator
	\begin{equation}\label{operador}
		T_\alpha f(x)=\int_{\mathbb{R}^{n}}K(x,y) f(y)dy,
	\end{equation}
	with $K(x,y)=k_1( x-A_1y)\cdots k_m( x-A_my)$,
	where $A_{i}$ are certain invertible matrices for all $i=1,\ldots,m$,
	and $f\in L^\infty_{loc}(\mathbb{R}^n)$.
	
	In the case  $A_i=a_iI$ with $a_i\in \mathbb{R}$ for all $i=1,\ldots, m$,  T.\@ Godoy and M.\@ Urciuolo in \cite{GU2}  obtained  the $L^p(\mathbb{R}^{n},dx)-L^q(\mathbb{R}^{n}, dx)$ boundedness of this operator for $0\leq \alpha < n$, $1 < p <\frac n\alpha$ and $\frac 1q=\frac 1p-\frac \alpha n$. In \cite{RoU}, for  $\Omega _i$ being smooth functions, P.\@ Rocha and M.\@ Urciuolo considered the operator $T_\alpha$ for  matrices $A_1,\dots, A_m$ satisfying the following hypothesis
	
	\
	
	
	(H)\ \ \emph {$A_{i}$  is invertible and  $A_i-A_j$ is invertible for $i\not = j$, $1\leq i, j\leq m$.}
	
	
	\
	
	
	\noindent They obtain that $T_\alpha$ is  bounded from $H^{p}(\mathbb{R}^{n},dx)$ into $L^{q}(\mathbb{R}%
	^{n},dx),$ for $0<p<\frac{n}{\alpha}$ and $\frac{1}{q}=\frac{1}{p}-\frac{\alpha }{n}$.  
	In this paper, they also show that $H^p-H^q$ cannot be expected. This is an important difference with respect to the case $m=1$.
	
	
	The weighted $L^p(\mathbb{R}^{n},dx)-L^q(\mathbb{R}^{n},dx)$ boundedness of $T_\alpha$ was obtained by M.\@ S.\@ Riveros and M.\@ Urciuolo in \cite{RiU2}.
	In this paper, we will prove a general version of this result for kernels satisfying certain conditions, which we define below.
	
	We denote by  $|x|\sim R$ the set $\{x\in \mathbb{R}^n: R<|x|\leq 2R\}$  and for $1\leq r\leq \infty$
	$$||f||_{r, |x|\sim R}=%||f\chi_\{|x|\sim R\}||_{r, B(0,2R)}
	\left(\frac 1{|B(0,2R)|}\int _{B(0,2R)}|f|^r\chi_{|x|\sim R}\right)^{\frac 1r}.$$
	In \cite {BLRi}, the authors introduced the following definition.
	
	\begin{definition}
		Given  $0\leq \alpha<n$ and  $1\leq r\leq \infty$ we say that
		$k \in H_{r, \alpha}$ if there exist $c\geq 1$ and  $C>0$ such that for all  $y\in \mathbb{R}^n$ and  $R>c|y|$
		$$
		\sum_{m=1}^{\infty}(2^mR)^{n-\alpha}||k(.-y)-k(.)||_{r, |x|\sim 2^mR}\leq C.
		$$
	\end{definition} 
	
	
	We consider the operator $T_\alpha$  defined in (\ref{operador}) where, for  $1\leq i\leq m$, $k_i$ is given by (\ref{ki}) and the matrices $A_i$ satisfy the hypothesis (H). Moreover, the operator also satisfies the following two conditions:
	
	\
	
	($\mathrm{H}_1$)\ \ There exists $\{p_i\}_{i=1}^{m}$ such that  $p_i>q_i$, $\Omega _i\in L^{p_i}(\Sigma)$ for all $1 \leq i \leq m$ and $\sum_{i=1}^{m} \frac{1}{p_i} < 1$,
	
	\
	
	
	($\mathrm{H}_2$)\ \  $k_i\in H_{p_i, \frac n{q_i' }-\alpha}$ for all $1 \leq i \leq m$.
	
	
	\medskip
	
	Let  $t\geq 1$  be  defined by  $\frac 1{p_1}+\dots+\frac 1{p_m}+\frac 1 t=1$ and let us consider the following useful coefficient $a=a(t,q)=\frac{tq}{t+q-tq}$.
	
	\ 
	
	This paper aims to prove the following result.
	
	\begin{theorem}\label{T-HpLq} Let $0< \alpha<n$, $0 < p <1$, $\frac 1q=\frac 1p-\frac {\alpha}n$ and let $T_\alpha$ be the integral operator defined by (\ref{operador}). We suppose that for $1\leq i\leq m$, the matrices $A_i$ and the functions $\Omega_i$ satisfy the hypotheses ($H$), ($H_1$) and ($H_2$). Let  $0<s\leq p$ and let $w$ be a weight such that $w^{\max\{a(t,q), \frac{n}{(n-\alpha)s}\}} \in \mathcal{A}_1$ and $w$ satisfying $w(A_ix)\leq Cw(x)$ $a.e.x \in \mathbb{R}^n$ for all $1\leq i\leq m$. Then $T_{\alpha}$ can be extended to an $H^{p}_{w^p}(\mathbb{R}^{n})-L^{q}_{w^q}(\mathbb{R}^{n})$ bounded operator.
	\end{theorem}	
	
	\
	The paper is organized as follows. In Section 2, we introduce the necessary definitions and preliminary results: weighted theory and Hardy spaces. 
	In Section 3, we prove an extension theorem \ref{Extension} and the main result.
	
	\section{Preliminaries}
	
	\subsection{Weighted Theory}
	A weight is a function $w\in L^1_{loc}(\mathbb{R}^n, dx)$ such that $w$ takes values in $(0,\infty)$ almost everywhere.
	Given a weight $w$ and $0<p<\infty$, we define $L^{p}_w(\mathbb{R}^{n})$ the space of all functions $f:\mathbb{R}^{n}\rightarrow \mathbb{R}$ that satisfy $||f||^{p}_{L^{p}_w}:=\int_{\mathbb{R}^{n}} |f(x)|^{p}w(x)dx< \infty$. When $p=\infty$, we have $L^{\infty}_w(\mathbb{R}^{n}):=L^{\infty}(\mathbb{R}^{n})$ and $||f||_{L^{\infty}_w}=||f||_{L^{\infty}}$. 
	
	Let $f \in L^{1}_{loc}(\mathbb{R}^{n})$. Recall that the Hardy-Littlewood maximal of $f$ is
	$$Mf(x)=\sup_{x \in B} \frac{1}{|B|} \int_B |f(y)| dy,$$
	where the supremum is taken over all balls $B$ containing $x$. 
	
	The Muckenhoupt class $\mathcal{A}_{p}$, $1<p<\infty $, is defined as the class of weights $w $ such that
	
	\begin{equation*}
		[w]_{\mathcal{A}_p}:=\sup_{B} \left[ \left( \frac{1}{\lvert B \rvert} \displaystyle \int_B w
		\right) \left( \frac{1}{\lvert B \rvert} \displaystyle \int_B w^{-\frac{
				1}{p-1}} \right)^{p-1} \right] < \infty,
	\end{equation*}
	where $B$ is a ball in $\mathbb{R}^n$.
	
	For $p=1$, $\mathcal{A}_1$ is the class of weights $w$ satisfying that
	there exists $c>0$ such that
	
	\begin{equation*}
		Mw(x) \leq c w(x) \ a.e. \ x \in \mathbb{R}^n.
	\end{equation*}
	
	We denote by $\left[ w \right]_{\mathcal{A}_1}$ the infimum of the
	constant $c$ such that $w$ satisfies the above inequality.
	Equivalently, we have $w \in \mathcal{A}_1$ if there exists $C>0$ such that for all ball $B$ we have
	$$\frac{1}{|B|} \int_B w(x) dx \leq C essinf_{x \in B} w(x).$$
	
	\begin{remark}\label{Remark A1}
		Notice that if $w \in \mathcal{A}_1$, then $w^{r} \in \mathcal{A}_1$ for $0<r<1$.
	\end{remark}
	
	%It is well known (see \cite{MW}) that if $w$ is a weight then $M_{\alpha }$ is bounded from $L^{p}\left(\mathbb{R}^{n},
	%w^{p}\right) $ into $L^{q}\left(\mathbb{R}^{n}, w^{q}\right) ,$ for $1<p<\frac{n}{\alpha }$
	%and $\frac{1}{q}=\frac{1}{p}-\frac{\alpha }{n},$ if and only if
	%\begin{equation}
	%\sup_{B}\left[\left( \frac{1}{\left\vert B\right\vert }\int_{B}w^{q}\right)
	%^{\frac{1}{q}}\left( \frac{1}{\left\vert B\right\vert }\int_{B}w^{-p^{\prime
			%}}\right) ^{\frac{1}{p^{\prime }}}\right]<\infty , \label{Apq}
	%\end{equation}%
	%where $\frac 1p+\frac 1{p'}=1$. The class of weights that satisfy (\ref{Apq}) is called $\mathcal{A}(p,q)$.
	
	%{\color{red} Chequear si verdaderamente usamos la clase $A(p,q)$}.
	
	Throughout this paper we understand that for $p=\infty$,  $(\int_E|f|^p)^{\frac 1p}$ stands for  $||f\chi_E||_\infty$, for any $E$ is a measurable set .
	%we define  the class $\mathcal{A}(p,q)$  still by  (\ref{Apq}) for all $1\leq p\leq \infty$ and $1\leq q\leq \infty$.
	For $p\geq 1 $, $\mathcal{A}_p$ denotes the classical Muckenhoupt class of weights. 
	%We note that $w\in \mathcal{A}(p,q)$ if only if $w^q\in \mathcal{A}_{1+q/p'}$ if only if $w^{-p'}\in \mathcal{A}_{1+p'/q}$. As a particular case $w\in \mathcal{A}(p,p)$ is equivalent  to $w^p \in \mathcal{A}_p$.  
	
	We recall that  $ \mathcal{A}_\infty=\cup_{p\geq 1} \mathcal{A}_p$, then, if $w \in \mathcal{A}_p$, with $1 < p < \infty$, then there exists $1<q<p$ such that $w \in \mathcal{A}_q$. We denote by $\tilde{q}_w=\inf\{q>1 : w \in \mathcal{A}_q\}$ the critical index of $w$. %We also recall that the statement $w\in \mathcal{A}(\infty,\infty)$ is equivalent to $w^{-1} \in \mathcal{A}_1$ 
	
	\bigskip
	
	In this paper, we assume that the weight $w$ satisfies the additional hypothesis: \noindent \emph{There exists $c>0$  such that}
	\begin{equation}\label{hipotesis w}
		w(A_ix)\leq cw(x),
	\end{equation}
	\emph{  a.e.} $x\in\mathbb{R} ^n $, $1\leq i\leq m$.
	
	A weight $w$ satisfies the reverse H\"{o}lder's inequality with exponents $s>1$, denoted by $w \in RH_s$ if there exists $C>0$ such that for every ball $B$ we have that 
	\begin{equation}\label{cond s}
		\Big( \frac{1}{|B|} \int_B w(x)^{s} dx \Big)^{\frac{1}{s}} \leq C\frac{1}{|B|} \int_B w(x)dx.
	\end{equation}
	The best possible constant is denoted by $[w]_{RH_s}$. We observe that if $w \in RH_s$, then by H\"{o}lder's inequality, $w \in RH_t$ for all $1 < t < s$ and $[w]_{RH_t} \leq [w]_{RH_s}$.
	Moreover, if $w \in RH_s$, $s>1$, then $w \in RH_{s+\epsilon}$ for some $\epsilon >0$. We denote by $r_w=\sup\{r>1 : w \in RH_r\}$ the critical index of $w$ for the reverse H\"{o}lder's condition.
	It is well known that $w \in \mathcal{A}_{\infty}$ if and only if $w \in RH_s$ for some $s>1$. Thus, $1<r_w \leq \infty$ for all $w \in \mathcal{A}_{\infty}$.\newline
	
	Given a weight $w$, $0<p<\infty$ and a measurable set $E$ we set 
	$$w^p(E)=\int_E w(x)^p dx.$$ 
	
	\begin{lemma}
		If $w \in \mathcal{A}_p$ for some $1\leq p< \infty$, then the measure $w(x)dx$ is doubling: for all $\lambda >1$ and all ball $B$ we have
		\begin{equation*}
			w(\lambda B) \leq \lambda^{np} [w]_{\mathcal{A}_p} w(B),
		\end{equation*}
		where $\lambda B$ denotes the ball with the same center as $B$ and radius $\lambda$ times the radius of $B$.
	\end{lemma}
	The following result is an immediate consequence of the reverse H\"{o}lder's condition.
	
	\begin{lemma}\label{weightRH}
		For $0<\alpha<n$, let $0<p<\frac{n}{\alpha}$ and $\frac{1}{q}=\frac{1}{p}-\frac{\alpha}{n}$. If $w^p  \in RH_{\frac{q}{p}}$ then
		\begin{equation}
			[w^p(B)]^{-\frac{1}{p}} [w^q(B)]^{\frac{1}{q}}\leq [w^p]^{1/p}_{RH_{q/p}} |B|^{-\frac{\alpha}{n}},
		\end{equation}
		for each ball $B$ in $\mathbb{R}^n$.
	\end{lemma}
	
	\
	
	
	
	\subsection{Weighted Hardy Spaces}
	
	Topologize $\mathcal{S}(\mathbb{R}^n)$ by the collection of semi-norms
	$\|\cdot\|_{\alpha,\beta}$, with $\alpha$ and $\beta$ multi-indices, given by
	\[
	\|\phi\|_{\alpha,\beta}
	=
	\sup_{x\in\mathbb{R}^n}
	\left|x^\alpha \partial^\beta \phi(x)\right|.
	\]
	
	For each $N\in\mathbb{N}$, we set
	\[
	S_N
	=
	\left\{
	\phi\in\mathcal{S}(\mathbb{R}^n):
	\|\phi\|_{\alpha,\beta}\le 1,\,
	|\alpha|,\,|\beta|\le N
	\right\}.
	\]
	
	Let $f\in\mathcal{S}'(\mathbb{R}^n)$, we denote by $M_N$ the grand maximal operator given by
	\[
	M_Nf(x)
	=
	\sup_{t>0}
	\sup_{\phi\in S_N}
	\left|
	t^{-n}\phi(t^{-1}\cdot)*f(x)
	\right|.
	\]
	
	Given a weight $w\in \mathcal{A}_\infty$ and $p>0$, the weighted Hardy space
	$H_w^p(\mathbb{R}^n)$ consists of all tempered distributions $f$ such that
	\[
	\|f\|_{H_w^p(\mathbb{R}^n)}
	=
	\|M_Nf\|_{L_w^p(\mathbb{R}^n)}
	=
	\left(
	\int_{\mathbb{R}^n}
	[M_Nf(x)]^p
	w(x)\,dx
	\right)^{1/p}
	<\infty.
	\]
	
	Let $\phi\in\mathcal{S}(\mathbb{R}^n)$ be a function such that
	\[
	\int_{\mathbb{R}^n}\phi(x)\,dx\neq 0.
	\]
	For $f\in\mathcal{S}'(\mathbb{R}^n)$, we define the maximal function
	$M_\phi f$ by
	\[
	M_\phi f(x)
	=
	\sup_{t>0}
	\left|
	t^{-n}\phi(t^{-1}\cdot)*f(x)
	\right|.
	\]
	
	For $N$ sufficiently large, we have
	\[
	\|M_\phi f\|_{L_w^p}
	\simeq
	\|M_Nf\|_{L_w^p},
	\qquad
	\text{(see \cite{ST}).}
	\]
	
	In the sequel we consider the following set
	\[
	\mathcal{D}_{0}
	=
	\left\{
	\phi\in\mathcal{S}(\mathbb{R}^n):
	\widehat{\phi}\in C_c^\infty(\mathbb{R}^n)
	\text{ and }
	0\notin\operatorname{supp}(\widehat{\phi})
	\right\}.
	\]
	
	Here, we recall some results that are useful in the following chapters.
	
	\begin{theorem}[\cite{ST}]
		Let $w$ be a doubling weight on $\mathbb{R}^{n}$. Then, $\hat{\mathcal{D}_0}$ is dense in $H^{p}_w(\mathbb{R}^{n})$, $0<p<\infty$.
	\end{theorem}
	
	We recall the definition of $w$-$(p,p_0,d)$ atom given in \cite{Rocha}.
	
	Let $w \in \mathcal{A}_{\infty}$ with critical index $\tilde{q}_w$ and critical index $r_w$ for the reverse H\"{o}lder's condition. Let $0<p\leq 1$, $\max\{1,p\frac{r_w}{r_w-1}\} <p_0 \leq \infty$, and $d \in \mathbb{Z}$ such that $d \geq \lfloor n ( \frac{\tilde{q}_w}{p}-1)\rfloor $ we say that a function $a(\cdot)$ is a $w$-$(p,p_0,d)$ atom centered in $x_0 \in \mathbb{R}^{n}$ if
	
	$$\begin{aligned}
		(a1)  \quad &  a \in L^{p_0}(\mathbb{R}^{n}) \text{ with support in the ball $B=B(x_0,r)$.} \\
		(a2) \quad  &  ||a||_{L^{p_0}(\mathbb{R}^{n})} \leq |B|^{\frac{1}{p_0}}w(B)^{-\frac{1}{p}}. \\
		(a3) \quad & \int_{\mathbb{R}^{n}} x^{\alpha} a(x) dx=0 \text{ for all multi-index $\alpha$ such that $|\alpha|\leq d$}.
	\end{aligned}
	$$
	
	\begin{lemma}[\cite{Rocha}]\label{Lemma8} Let $w \in \mathcal{A}_{\infty}$ with critical index $\tilde{q}_w$ and critical index $r_w$ for the reverse H\"{o}lder's condition. If $a(\cdot)$ is a $w$-$(p,p_0,d)$ atom, then $a(\cdot) \in H^{p}_w(\mathbb{R}^{n})$. Moreover, there exists a positive constant $C$ independent of the atom $a$ such that $||a||_{H^{p}_w} \leq C$. 
	\end{lemma}
	\begin{theorem}[\cite{Rocha}]\label{Theorem9}
		Let $f \in \hat{\mathcal{D}_0}$, and $0<p \leq 1$. If $w \in \mathcal{A}_{\infty}$, then there exist a sequence of $w$-$(p,p_0,d)$ atoms $\{a_j\}$ and a sequence of scalar $\{\lambda_j\}$ with $\sum_j |\lambda_j|^{p} \leq c||f||^{p}_{H^{p}_w}$ such that $f=\sum_j \lambda_ja_j$, where the convergence is both in $L^{s}(\mathbb{R}^{n})$ and pointwise, for each $1<s<\infty$.
	\end{theorem}
	
	\ 
	
	
	
	\section{Proof of Main results}
	
	In order to prove our main result, we first prove an extension result and a lemma about the uniform boundedness of the operator $T_{\alpha}$. 
	
	\begin{theorem}\label{Extension}
		Let $T$ be a bounded linear operator from $L^{p_0}(\mathbb{R}^{n})$ to $L^{q_0}(\mathbb{R}^{n})$ for some $1<p_0<\infty$ and $p_0<q_0<\infty$. Let $w \in \mathcal{A}_{\infty}$ with critical index $r_w$, $0<p \leq \min \{1, \frac{r_w-1}{r_w} p_0 \}$, and let $q\geq p$ be a real number. Then $T$ can be extended to a bounded linear operator $H^{p}_w(\mathbb{R}^{n})-L^{q}_w(\mathbb{R}^{n})$ if and only if $Ta$ is uniformly bounded into the $L^{q}_w$ norm for all $w-(p,p_0,d)$ atoms $a$.
	\end{theorem}
	
	\begin{proof}
		By Lemma \ref{Lemma8}, we have that $a \in H^{p}_w(\mathbb{R}^{n})$ for all $w$-$(p,p_0,d)$ atom. Then, if $T$ can be extended to an $H^{p}_w(\mathbb{R}^{n})-L^{q}_w(\mathbb{R}^{n})$, we obtain that $||Ta||_{L^{q}_w} \leq c_p||a||_{H^{p}_w}$. Therefore, by Lemma \ref{Lemma8} we have that $||Ta||_{L^{q}_w} \leq c_p$ for all $a$ $w$-$(p.p_0,d)$ atom.
		
		Conversely, by Theorem \ref{Theorem9}, given $f \in \hat{\mathcal{D}_0}$ there exist a sequence of $w$-$(p,p_0,d)$ atoms $\{a_j\}$ and a sequence of scalar $\{\lambda_j\}$ with $\sum_j |\lambda_j|^{p} \leq c||f||^{p}_{H^{p}_w}$ such that $f=\sum_j \lambda_ja_j$, where the convergence is both in $L^{p_0}(\mathbb{R}^{n})$ and pointwise. 
		Since $T$ is a bounded operator from $L^{p_0}(\mathbb{R}^{n})$ in $L^{q_0}(\mathbb{R}^{n})$, we have that $\sum_j \lambda_jTa_j$ converges to $Tf$ in $L^{q_0}(\mathbb{R}^{n})$.
		Then, there exists a subsequence $\{r_n\}$ such that $\lim_{n \to \infty} \sum_{j=1}^{r_n} \lambda_j Ta_j(x)=Tf(x)$ a.e.$x \in \mathbb{R}^{n}$.
		Then, 
		\begin{equation}\label{Tpuntual}
			|Tf(x)| \leq \sum_j |\lambda_jTa_j(x)|, \quad a.e.x \in \mathbb{R}^{n}.
		\end{equation}
		
		In the follow we analyze the cases $p\leq q<1$ and $q\geq 1$. First, if $p\leq q<1$, \eqref{Tpuntual} implies that 
		$$ \int_{\mathbb{R}^{n}} |Tf(x)|^{q} w(x) dx \leq \sum_j |\lambda_j|^{q} \int_{\mathbb{R}^{n}} |Ta_j(x)|^{q} w(x) dx.$$
		If we suppose that $||Ta||_{L^{q}_w} \leq C_p$ for all $w$-$(p.p_0,d)$ atom $a$, we have that 
		$$ \Big(\int_{\mathbb{R}^{n}} |Tf(x)|^{q} w(x) dx \Big)^{\frac{1}{q}} \leq C_p \Big(\sum_j |\lambda_j|^{q}\Big)^{\frac{1}{q}}.$$
		It is easy see that $ \left( \sum_j |\lambda_j|^q \right)^{\frac{1}{q}} \leq \Big( \sum_j|\lambda_j|^{p}\Big)^{\frac{1}{p}} \leq c ||f||_{H^{p}_w}$, 
		so we have that 
		$$ \Big(\int_{\mathbb{R}^{n}} |Tf(x)|^{q} w(x) dx \Big)^{\frac{1}{q}} \leq C_p \Big(\sum_j |\lambda_j|^{q}\Big)^{\frac{1}{q}}  \leq C_p ||f||_{H^{p}_w},$$
		
		for all $f \in \hat{\mathcal{D}_0}$, and the Theorem follows by the density of $\hat{\mathcal{D}_0}$ in $H^{p}_w(\mathbb{R}^{n})$.
		
		Now, if $q\geq 1$, by Minkowsky's inequality, monotone convergence theorem and \eqref{Tpuntual} implies that 
		$$
		\begin{aligned}
			\left(\int_{\mathbb{R}^{n}} |Tf(x)|^{q} w(x) dx\right)^{\frac{1}{q}} &\leq \left(\int_{\mathbb{R}^{n}} \left(\sum_{j=1}^{\infty} |\lambda_j| |Ta_j(x)|\right)^{q} w(x) dx\right)^{\frac{1}{q}}\\
			&\leq  \lim_{n \to \infty} \left(\int_{\mathbb{R}^{n}} \left(\sum_{j=1}^{n} |\lambda_j| |Ta_j(x)|\right)^{q} w(x) dx\right)^{\frac{1}{q}}\\
			&\leq  \lim_{n \to \infty} \Big\| \sum_{j=1}^{n} |\lambda_j| |Ta_j| \Big\|_{L^{q}(w)}\\
			&\leq \lim_{n \to \infty} \sum_{j=1}^{n}|\lambda_j| \ || Ta_j||_{L^{q}(w)}\\
			&\leq C \lim_{n \to \infty} \sum_{j=1}^{n}|\lambda_j|.
		\end{aligned}
		$$
		
		Since $0<p\leq 1$ , we have that $\sum_j |\lambda_j| \leq \Big( \sum_j|\lambda_j|^{p}\Big)^{\frac{1}{p}} \leq c ||f||_{H^{p}_w}$, and then  we have that 
		$$ \Big(\int_{\mathbb{R}^{n}} |Tf(x)|^{q} w(x) dx \Big)^{\frac{1}{q}} \leq C_p ||f||_{H^{p}_w},$$
		for all $f \in \hat{\mathcal{D}_0}$, and the Theorem follows by the density of $\hat{\mathcal{D}_0}$ in $H^{p}_w(\mathbb{R}^{n})$.
	\end{proof}
	
	\medskip
	
	Next, under suitable assumptions, we prove that the operator defined by \eqref{operador}, when applied to a $w^{p}$-$(p,p_0,d)$ atom, is uniformly bounded in the corresponding $L^q$ space. 
	
	\begin{lemma}\label{TL}
		Let $0< \alpha<n$, $0 < p \leq 1$, $0< q$ satisfies $\frac 1q=\frac 1p-\frac {\alpha}n$ and let $T_\alpha$ be the integral operator defined by \eqref{operador}. We suppose that for $1\leq i\leq m$, the matrices $A_i$ and the functions $\Omega_i$ satisfy the hypothesis ($H$), ($H_1$) and ($H_2$). Suppose that  $w^{\max\{a(t,q), p\}} \in \mathcal{A}_1$ with $w$ satisfying \eqref{hipotesis w}, then there is $C>0$ independent of $a$ such that
		$$
		\left(\int_{\mathbb{R}^{n}} |T_\alpha a(y)|^qw^q(y)dy\right)^{1/q} \leq C  
		$$
		for all $w^{p}$-$(p,p_0,d)$ atom $a$ with $supp(a) \subseteq B(x_0,R_0)$, where  $d \in \mathbb{Z}$ such that $d \geq \lfloor n ( \frac{\tilde{q}_w}{p}-1)\rfloor $ and  $p_0> \max\{1,\frac{q}{q-1}(1-\delta_{q,1})\}$, with $\delta_{q,1}=1$ if $q=1$ and $0$ otherwise.
	\end{lemma}
	
	
	\begin{proof}
		We consider a $w^p-(p,p_0, d)$ atom
		$a$  with $supp(a) \subseteq B(x_0,R_0)$ such that $R_0>\max\{|x_0|,1\}$. Since $r_w \geq 1$, then $d\geq n(\frac 1p-1)$. For all $1 \leq i \leq m$, let $B_M=B(x_0, 2M\sqrt n  R_0)$, where $
		M=\max_{1\leq i\leq m}\{|A_i|\}$. 
		
		We have that
		$$
		\begin{aligned}
			||T_\alpha a||_{L^q(w^q)}&=\left(\int_{\mathbb{R}^n} |T_\alpha a(y)|^qw^q(y)dy\right)^{1/q}\\
			&\leq   C \left( \sum_{i=1}^m\left(\int_{A_i(B_M)} |T_\alpha a(y)|^qw^q(y)dy\right)^{1/q}+\left(\int_{(\cup_i A_i(B_M))^c} |T_\alpha a(y)|^qw^q(y)dy\right)^{1/q}\right)\\
			&:=\sum _{i=1}^mI_i+I_{m+1}.
		\end{aligned}
		$$
		
		Note that, since $0<p\leq 1$ and $\frac{1}{q}=\frac{1}{p} - \frac{\alpha}{n}$, if $p \leq \frac{n}{n+\alpha}$ then $q\leq 1$ and if $p> \frac{n}{n+\alpha}$ then $q>1$; the constant $C$ take two values: $C=1$ when $q>1$ and $C=2^{\frac{1}{q} -1}$. Furthermore, when $q \geq 1$, the preceding inequality is the triangle inequality in $L^q$, and when $q<1$, this inequality follows from Jensen's inequality for concave functions.
		
		Since $p_0 > \frac{q-1}{q}$, we have that $\frac{q_0}{q}>1$. 
		To estimate $I_i$, we apply the H\"{o}lder's inequality with $\frac{q_0}{q}$, then we use that $w^{q} \in RH_{(\frac{q_0}{q})'}$ and the fact that $w(A_ix) \leq C w(x), \ a.e.x\in \mathbb{R}^n$. Then,
		
		$$
		\begin{aligned}
			I_i & =\int_{A_i(B_M)} |T_\alpha a(y)|^qw^q(y)dy \\
			&\leq ||T_{\alpha}a||^{q}_{L^{q_0}} \Big(\int_{A_i(B_M)}[w^{q}(x)]^{(\frac{q_0}{q})'} dx\Big)^{\frac{1}{(\frac{q_0}{q})'}}\\
			&\leq ||T_{\alpha}a||^{q}_{L^{q_0}} \Big(|det(A_i)|^{n}\int_{B_M}[w^{q}(A_iy)]^{(\frac{q_0}{q})'} dy\Big)^{\frac{1}{(\frac{q_0}{q})'}}\\
			& \leq M^{n(\frac{q_0}{q})'}||T_{\alpha}a||^{q}_{L^{q_0}} \int_{B_M} c \ [w^{q}(y)]^{(\frac{q_0}{q})'} dy\Big)^{\frac{1}{(\frac{q_0}{q})'}}\\
			& \leq C M^{n(\frac{q_0}{q})'}||a||^{q}_{L^{p_0}} \int_{B_M}  \ [w^{q}(y)]^{(\frac{q_0}{q})'} dy\Big)^{\frac{1}{(\frac{q_0}{q})'}}\\
			& \leq C|B(x_0,R_0)|^{\frac{q}{p_0}} (w^{p}(B(x_0,R_0))^{-\frac{q}{p}} |B_M|^{\frac{1}{(\frac{q_0}{q})'}} \Big(\frac{1}{|B_M|} \int_{B_M} w^{q}(x)dx\Big) \\
			& \leq C|B_{M}|^{\frac{q}{p_0}} (w^{p}(B_M)^{-\frac{q}{p}} |B_M|^{\frac{1}{(\frac{q_0}{q})'}} \Big(\frac{1}{|B_M|} \int_{B_M} w^{q}(x)dx\Big) \\
			& \leq C |B_M|^{\frac{q\alpha}{n}} (w^{p}(B_M))^{-\frac{q}{p}} w^{q}(B_M),
		\end{aligned}
		$$
		and so by Lemma \ref{weightRH} we obtain $I_i \leq C$.
		
		
		
		We now proceed to bound $I_{m+1}$.  We  first observe that
		$(\bigcup A_i(B_M))^{C} \subseteq B(0,2RM)^{C}$ with $R=\frac{2M \sqrt{n}R_0-|x_0|}{2} > 0 $.
		Indeed, since the matrices $A_i$ are invertibles for all $1 \leq i \le m$, if $y \notin A_i(B_M)$, we have that $y=A_i(v_i)$ where $v_i \notin B_M$. Then, 
		\begin{eqnarray*}
			2M \sqrt{n}R_0 &\leq &|v_i-x_0|\\
			&=&|A_i^{-1}(y)-A_i^{-1}(A_i(x_0))| \\
			&\leq &|A_i|^{-1} |y-A_i(x_0)|.
		\end{eqnarray*}
		Thus, 
		$$2M\sqrt{n}R_0 \leq \frac{1}{|A_i|} |y-A_i(x_0)|$$ 
		and hence  
		$$|A_i| 2M\sqrt{n}R_0 \leq |y-A_i(x_0)| \leq |y|+|A_i(x_0)| \leq |y| + |A_i| |x_0|.$$ 
		Then, $|A_i| (2M\sqrt{n}R_0-|x_0|) \leq |y|$, and since it is true for all  $1 \leq i  \leq m$ we obtain
		$$M(2M\sqrt{n}R_0-|x_0|) \leq |y|.$$
		
		Finally, we have that $(\cup A_i(B_M))^{C} \subseteq B(0,2RM)^{C}$ with $R=\frac{2M\sqrt{n}R_0-|x_0|}{2} > 0 $. On the other hand, if $|z-x_0|< R_0$, since $|z|-|x_0| <|z-x_0|$, we have that  
		$$ |z|  <  R_0 + |x_Q| < 2 R_0 = \frac{R}{\sqrt{n}M}.$$
		Therefore, we obtain $B(x_0, R_0) \subseteq B(0,\frac{R}{\sqrt{n}M})$. Thus,
		
		$$
		\begin{aligned}
			I_{m+1} &=  \left(\int_{(\cup A_i(B_M))^{C}} |T_\alpha a(y)|^q w^q(y)dy\right)^{\frac{1}{q}}\\
			&\leq \left(\int_{|y|> 2MR} |T_\alpha a(y)|^q w^q(y)dy\right)^{\frac{1}{q}}\\
			&=  \left(\int_{|y|> 2MR} \left |\int_{|z-x_0|< R_0} K(y,z) a(z) dz \right |^{q}
			w^q(y)dy\right)^{\frac{1}{q}}\\
			&=  \left(\int_{|y|> 2MR} \left |\int_{|z-x_0|< R_0} \left[ K(y,z) a(z)-K(y,0)a(z)\right] dz \right |^{q}
			w^q(y)dy\right)^{\frac{1}{q}}\\
			&\lesssim  \int _{|z|\leq \frac{R}{\sqrt{n}M}}|a(z)|\sum_{j=2}^\infty\left(\int_{|y|\sim M R2^{j}} |K(y,z)-K(y,0)|^q w^q(y) dy\right)^{\frac 1q}dz
		\end{aligned}
		$$
		
		
		We estimate now $\left| K(y,z)-K(y,0)\right| $ for $|y|>MR$,  and $|z|\leq  n R$.  It is easy to check that
		{\small		\begin{equation}\label{formki}
				\left| K(y,z)-K(y,0)\right| \leq\sum_{i=1}^m \left[
				\prod_{r=1}^i|k_{r-1}(y)||k_i(y-A_iz)-k_i(y)|\prod_{r=i}^m |k_{r+1}(y-A_rz)|\right]
			\end{equation}
		}		
		%		\begin{equation}\label{formki}
			%			\begin{aligned}
				%				&\qquad\qquad\qquad\qquad\left| K(y,z)-K(y,0)\right|\\
				%				&\leq\sum_{i=1}^m \left[
				%				\prod_{r=1}^i|k_{r-1}(y)||k_i(y-A_iz)-k_i(y)|\prod_{r=i}^m |k_{r+1}(y-A_rz)|\right],
				%			\end{aligned}
			%		\end{equation}
		where we define $k_{0}=k_{m+1}\equiv 1$. Now, by Jensen's inequality with $0 < q \leq 1$, and generalized Minkowsky inequality in the case $q>1$, we have that
		$$
		\begin{aligned}
			&\left(\int_{|y|\sim M R2^{j}} \left(\sum_{i=1}^m \left[
			\prod_{r=1}^i|k_{r-1}(y)||k_i(y-A_iz)-k_i(y)|\prod_{r=i}^m |k_{r+1}(y-A_rz)|\right]\right)^q w^q(y) dy\right)^{1/q}\\
			& \leq C  \sum_{i=1}^m\left(\int_{|y|\sim M R2^{j}} \left[
			\prod_{r=1}^i|k_{r-1}(y)|^q|k_i(y-A_iz)-k_i(y)|^q\prod_{r=i}^m |k_{r+1}(y-A_rz)|^q\right] w^q(y) dy\right)^{1/q}.
		\end{aligned}
		$$
		
		
		
		We have that $\sum_{i=1}^{m}\frac{1}{p_i}+\frac{1}{t}=1$, and then, $\sum_{r=1}^{m}\frac{q}{p_r} + \frac{t+q-tq}{t}=1$. 
		
		\
		
		We observe that since $\sum_{i=1}^{m}\frac{1}{p_i} < \sum_{i=1}^{m} \frac{1}{q_i} = 1 - \frac{\alpha}{n}$, we have that $\frac{1}{q}=\frac{1}{p}-\frac{\alpha}{n} > 1 -\frac{\alpha}{n} > \sum_{i=1}^{m}\frac{1}{p_i}$, and if $t<1$, then $0 < \sum_{i=1}^{m}\frac{q}{p_i}<1$.
		
		Using generalized H\"{o}lder´s inequality with exponents $\frac{p_1}{q},\frac{p_2}{q}\dots \frac{p_m}{q},\frac {t}{t+q-tq}$, we get
		$$		
		\begin{aligned}
			& \left(\int_{|y|\sim M R2^{j}}\prod_{r=1}^{i-1}|k_r(y)|^q \ |k_i(y-A_iz)-k_i(y)|^q\prod_{r=i+1}^m |k_{r}(y-A_rz)|^q w^q(y)dy\right)^{\frac 1q} \\
			& \leq \left(\prod_{r=1}^{i-1}|| k_r(\cdot)\chi_{|\cdot|\sim MR2^{j}}||_{p_r}||(k_i(\cdot-A_iz )-k_i(\cdot))\chi_{|\cdot|\sim M R2^{j}}||_{p_i}\right) \times \\
			&\left(\prod_{r=i+1}^m||k_{r}(\cdot-A_r z)\chi_{|\cdot|\sim M R2^{j}}||_ {p_r}
			|| w   \chi _{|\cdot|\sim M R2^{j}}||_{\frac {tq}{t+q-tq}}\right).
		\end{aligned}
		$$
		
		
		Observe that since  for $1 \leq  r\leq m$, $|A_rz| <M R $ and $|y|\sim   MR2^j $  then $MR 2^{j-1}<|y-A_rz|< MR2^{j+1}$.
		
		\ 
		
		Now, if $p_r<\infty$, since $\Omega_r$ is a homogeneous function of degree zero, we have that
		\begin{equation}
			\begin{aligned}
				||k_r(\cdot-A_rz)\chi _{|\cdot|\sim M R2^{j}}||_{p_r}&=
				\left(\int_{\{|y|\sim M R2^{j}\}} \frac{|\Omega_{r}(y-A_rz)|^{p_r}}{|y-A_rz|^{\frac{np_r}{q_r}}}dy\right)^{\frac 1{p_r}}\\
				&\leq C (R2^j)^{-\frac{n}{q_r}}\left(\int_{\{2^{j-1}MR<|y-A_lz|\leq 2^{j+1}MR\}}|\Omega_{r}(y-A_rz)|^{p_r}dy\right)^{\frac 1{p_r}}\\
				%&\leq C (R2^j)^{-\frac{n}{q_l}}\left(\int_{\{1<|u|\leq 8\}}|\Omega_{l}(2^jRu)|^{p_l}\,2^{jn}R^ndu\right)^{\frac 1{p_l}}\\
				&\leq C (2^jR)^{-\frac{n}{q_r}+\frac n{p_r}}\left(\int_{\{\frac M2<|u|\leq 2M\}}|\Omega_{r}(u)|^{p_r}du\right)^{\frac 1{p_r}}\\
				&\leq C (2^jR)^{-\frac{n}{q_r}+\frac n{p_r}}||\Omega_r||_{p_r}.
			\end{aligned}
		\end{equation}
		
		
		Analogously, we can prove that 
		$$||k_r(\cdot)\chi _{|\cdot|\sim M R2^{j}}||_{p_r} \leq C (2^jR)^{-\frac{n}{q_r}+\frac n{p_r}}||\Omega_r||_{p_r}.$$
		
		
		If $p_r=\infty$, we have that 
		\begin{equation}
			\begin{aligned}
				|k_r(\cdot-A_rz)\chi _{|\cdot|\sim M R2^{j}}|&=
				\frac{|\Omega_{r}(y-A_rz)|}{|y-A_rz|^{\frac{n}{q_r}}}\\
				&\leq C (R2^j)^{-\frac{n}{q_r}} |\Omega_{r}(y-A_rz)|\\
				&\leq C (R2^j)^{-\frac{n}{q_r}} ||\Omega_{r}||_{\infty}
			\end{aligned}
		\end{equation}
		Then, we have that 
		$$||k_r(\cdot-A_rz)\chi _{|\cdot|\sim M R2^{j}}||_{\infty} \leq C (2^jR)^{-\frac{n}{q_r}}||\Omega_r||_{\infty},$$
		and analogously, 
		$$||k_r(\cdot)\chi _{|\cdot|\sim M R2^{j}}||_{\infty} \leq C (2^jR)^{-\frac{n}{q_r}}||\Omega_r||_{\infty}.$$
		
		Also,
		$$
		\begin{aligned}
			& ||(k_i(\cdot-A_iz )-k_i(\cdot))\chi_{|\cdot|\sim M R2^{j}}||_{p_i}\\
			& \leq ( 2^jMR)^{\frac n{p_i}-\frac n{q_i}-\alpha}( 2^jMR)^{\frac n{q_i}-\frac n{p_i}+\alpha}||(k_i(\cdot-A_i z)-k_i(\cdot) )\chi_{|x|\sim 2^{j+1}MR}||_{p_i}.\\
		\end{aligned}
		$$
		\
		
		Since $w^{\frac {tq}{t+q-tq}}\in A_1 \subset A_{\infty}$, we have that $ w^{q} \in RH_{\frac{t}{t+q-tq}}$, and by the inequality \ref{cond s} we have that 
		$$\begin{aligned}
			||w \chi_{|\cdot|\sim M R2^{j}}||_{\frac {qt}{t+q-qt}}& \leq \left(\int_{B(0,MR2^{j})} w^{\frac{qt}{t+q-tq}}(x) dx \right)^{\frac{t+q-tq}{tq}}\\
			&\leq C |B(0,MR2^{j})|^{\frac{t+q-tq}{tq}} (\frac{1}{|B(0,MR2^{j})|})^{\frac{1}{q}} \left(\int_{B(0,MR2^{j})} w^{q}(x) dx\right)^{\frac{1}{q}}\\
			&=C |B(0,MR2^{j})|^{\frac{1}{t}-1} \left(\int_{B(0,MR2^{j})} w^{q}(x) dx \right)^{\frac{1}{q}} \\
			&= C (2^{j}R)^{\frac{n}{t}-n} \left(\int_{B(0,MR2^{j})} w^{q}(x) dx\right)^{\frac{1}{q}}.
		\end{aligned}$$
		
		Note that if $t=1$ then $p_i=\infty$ for all $i$. In this case, $\frac{t+q-tq}{t}=1$ and the above inequality is trivial since $w \in \mathcal{A}_1$.
		
		
		Since $k_i\in H_{p_i, \frac n{q_i'}-\alpha}$ and \eqref{q_i}, we have 
		$$\begin{aligned}
			\sum_{j=2}^\infty &\int_{|y|\sim M R2^{j}} \prod_{r=1}^{i-1}|k_r(y)|^{q}|k_i(y-A_iz)-k_i(y)|^q\prod_{r=i+1}^m |k_{r}(y-A_rz)|^q w^q(y)dy\\
			&\leq \left( CR^{\alpha-n}(w^{q}(B(0,2MR)))^{\frac 1q}\right) \times\\
			& \left(\sum_j 2^{j( \sum_r \frac n{p_r}-\sum_r \frac n{q_r}- \alpha+\frac{n}t-n } )(2^jR)^{\frac n{q_i}-\frac n{p_i} +\alpha} ||(k_i(\cdot-A_iz )-k_i(\cdot))\chi_{|\cdot|\sim M R2^{j}}||_{p_i}\right)\\
			&\leq CR^{\alpha-n}(w^{q}(B(0,2MR)))^{\frac 1q}\sum_j (2^j)^{-n}(2^jR)^{\frac n{q_i}-\frac n{p_i} +\alpha} ||(k_i(\cdot-A_iz )-k_i(\cdot))\chi_{|\cdot|\sim M R2^{j}}||_{p_i}\\
			&\leq CR^{\alpha-n}(w^{q}(B(0,2MR)))^{\frac 1q}\sum_j (2^jR)^{\frac n{q_i}-\frac n{p_i} +\alpha} ||(k_i(\cdot-A_iz )-k_i(\cdot))\chi_{|\cdot|\sim M R2^{j}}||_{p_i}\\
			&\leq  CR^{\alpha-n}(w^{q}(B(0,2MR)))^{\frac 1q}.
		\end{aligned}
		$$
		
		Then, since $B(x_0,R_0) \subseteq B(0,2R_0)$ we obtain that
		$$\begin{aligned}
			\sum_{i=1}^{m} \sum_{j=2}^\infty &\int_{|y|\sim M R2^{j}} \prod_{r=1}^{i-1}|k_r(y)|^{q}|k_i(y-A_iz)-k_i(y)|^q\prod_{r=i+1}^m |k_{r}(y-A_rz)|^q w^q(y)dy\\
			& \leq CR^{\alpha-n}(w^{q}(B(0,2MR)))^{\frac 1q}.
		\end{aligned}
		$$
		
		Then, we have that
		By \ref{TL}, we have that
		$$ \begin{aligned}
			I_{m+1}& \leq C   R^{\alpha-n}(w^{q}(B(0,2MR)))^{\frac 1q}\int _{|z|\leq\sqrt n R}|a(z)|\\
			&\leq C   R^{\alpha-n}(w^{q}(B(0,2MR)))^{\frac 1q} ||a||_{p_0} R^{\frac{n}{p'_0}}\\
			&\leq C  R^{\alpha-n}(w^{q}(B(0,2MR)))^{\frac 1q} |B(x_0,R_0)|^{\frac{1}{p_0}} w^{p}(B(x_0,R_0))^{-\frac{1}{p}} R^{\frac{n}{p'_0}}\\
			&\leq C  R^{\alpha-n}(w^{q}(B(0,2MR)))^{\frac 1q} |B(0,2R_0)|^{\frac{1}{p_0}} w^{p}(B(0,2R_0))^{-\frac{1}{p}} R^{\frac{n}{p'_0}}.
		\end{aligned}
		$$
		Sea $R_1:=\max\{R_0,2MR\} > 0$, with $R$ and $R_0$ as in Lemma \ref{TL}. Then, we have that
		\begin{equation*}
			I_{m+1} \leq C  R^{\alpha-n}(w^{q}(B(0,2R_1)))^{\frac 1q} |B(0,2R_1)|^{\frac{1}{p_0}} w^{p}(B(0,2R_1))^{-\frac{1}{p}} R_1^{\frac{n}{p'_0}}.
		\end{equation*}
		
		Lemma \ref{weightRH} and  $\frac{1}{q}=\frac{1}{p}-\frac{n}{\alpha}$, therefore, we have that
		
		$$w^p(B(0,2R_1))^{-\frac 1p}(w^{q}(B(0,2R_1)))^{\frac 1q} \leq CR_1^{-\alpha},
		$$
		and since $\frac{1}{p_0}+\frac{1}{p_0'}=1$ we obtain 
		
		$$I_{m+1}\leq C$$
		
		then we have proved that  
		
		$$
		||T_\alpha a||_{L^q(w^q)}\leq C
		$$
		for any $w^p(p,p_0, d)$-atom $a$. Then, we prove the result as desired.
		
	\end{proof}
	
	
	
	We are now in a position to prove our main result. Note that we are assuming that $w^{\max\{\frac{tq}{t+q-tq},\frac{n}{(n-\alpha)s}\}} \in \mathcal{A}_1$. An example of these weights in $\mathbb{R}^n$ is as follows.
	\begin{equation*}
		w(x)=|x|^{\frac{\beta}{\max\{\frac{tq}{t+q-tq},\frac{n}{(n-\alpha)s}\}}}
	\end{equation*}
	with $-n<\beta<0$.  Therefore, we may take orthogonal matrices $A_i$, and condition \eqref{hipotesis w} is then satisfied.
	
	
	\begin{proof}{(Theorem \ref{T-HpLq})}
		First, we observe that $w^{\max\{\frac{tq}{t+q-tq},\frac{n}{(n-\alpha)s}\}} \in \mathcal{A}_1$ implies $w^{\frac{tq}{t+q-tq}} \in \mathcal{A}_1$ and $w^{\frac{n}{(n-\alpha)s}} \in \mathcal{A}_1$ by Remark \ref{Remark A1}.
		
		For $s\leq p \leq 1<\frac{n}{n-\alpha}$, we have $\frac{n}{(n-\alpha)s}>\frac{n}{(n-\alpha)p}>\frac{n-\alpha}{(n-\alpha)p} = \frac{1}{p}$. Therefore, since  $w^{\frac{n}{(n-\alpha)s}} \in \mathcal{A}_1$  for all $s \leq p \leq 1$, and again by Remark \ref{Remark A1} we have that $w, w^{\frac{1}{p}}$, $w^{p}$ and $w^{q}$ belong to $\mathcal{A}_1$ for $p,q$ satisfying $\frac{1}{q}=\frac{1}{p}-\frac{\alpha}{n}$.
		Then, as $w^{\frac{1}{p}} \in \mathcal{A}_1 \subset \mathcal{A}_{\infty}$, we have $r_w>\frac{1}{p}$, (since $w^{r} \in \mathcal{A}_{\infty}$ if and only if $w \in RH_r$), and therefore $p>\frac{1}{r_w}$, which implies $\frac{1}{p}-1<\frac{r_w-1}{r_wp}$.
		
		Since $p>\frac{1}{r_w}$, we have that $-p<-\frac{1}{r_w} \quad 1-p <1-\frac{1}{r_w}=\frac{r_w-1}{r_w}$. We consider $1 < p_0$ such that  $\frac{1}{p_0} < \min\{\frac{1-p}{p},1\}$.
		In particular, we have that $\frac{1}{p_0} < \frac{1-p}{p} < \frac{r_w-1}{r_wp}$, and then $p < \frac{(r_w-1)p_0}{r_w}$. We consider $q_0$ such that $\frac{1}{q_0}=\frac{1}{p_0}-\frac{\alpha}{n}$. Since $\max\{1,\frac{q}{q-1}\} <p_0<\infty$, and $T_{\alpha}: L^{p_0}(\mathbb{R}^{n}) \rightarrow L^{q_0}(\mathbb{R}^{n})$ is a bounded operator (see \cite{FR}), by Theorem \ref{Extension}, it is enough to show that there is an absolute constant $C$ such that $||T_\alpha a||_{L^q(w^q)}\leq C$ for any $w^p(p,p_0, d)$-atom $a$ with center in $x_0 \in \mathbb{R}^{n}$. The theorem follows by choosing $p_0$ in such a way that the Lemma \ref{TL} can be applied.
		
		
		
	\end{proof}
	
	%\section{Technical Lemma}
	
	
	
	\ 
	
	\textbf{Acknowledgement.} We would like to express our gratitude to Rocha P.(Universidad Nacional del Sur) and Urciuolo M. (Universidad Nacional de Córdoba) for their helpful clarifications throughout the course of this work.
	
	
	\begin{thebibliography}{999}
		\parskip=0.1cm
		
		\bibitem{BLRi}  Bernardis A., Lorente M.,  Riveros M.S.,  Weighted Inequalities for fractional integral operators with kernel satisfying H\"{o}rmander type conditions,  \emph{ Mathematical Inequalities and Applications},
		\textbf{14} (4), 881-895, (2011).
		
		\bibitem{B} Bownik, Marcin. "Boundedness of operators on Hardy spaces via atomic decompositions." Proceedings of the American Mathematical Society 133.12 (2005): 3535-3542.
		
		\bibitem{C} Coifman, Ronald. "A real variable characterization of $ H^{p} $." Studia Mathematica 51.3 (1974): 269-274.
		
		\bibitem {DL} Ding Y., Lu S., Weighted norm inequalities for fractional integral operators with rough kernel, \emph{Can. J. Math.} \textbf{50} (1), 29-39, (1998).
		
		\bibitem{DL2} Ding Y., Lu S., Homogeneous fractional integrals on Hardy spaces,  \emph{T\^{o}hoku Math. J.}  \textbf{52} 153-162,(2000).
		
		\bibitem{DLeLi} Ding Y.,  Lee M., Lin C.,  Fractional integrals on weighted Hardy spaces,  \emph{J.Math. Anal.  Appl.} \textbf{282}, 356-368 (2003).
		
		\bibitem{D} Duoandikoetxea J., An\'{a}lisis de Fourier. Ediciones de la
		Universidad Aut\'{o}noma de Madrid, Editorial Siglo XXI, (1990).
		
		\bibitem{Du} Duoandikoetxea J., Weighted norm inequalities for homogeneous singular integrals, \emph{Trans. Amer. Math. Soc} \textbf{336}, 869-880, (1993).
		
		\bibitem{FS} Fefferman, Charles, and Elias M. Stein. "$H^p$ spaces of several variables." (1972): 137-193.
		
		\bibitem{FR} Firnkorn, Gonzalo H. Ibañez, and María Silvina Riveros. "Certain fractional type operators with Hörmander conditions." Annales Fennici Mathematici 43.2 (2018): 913-929.
		
		\bibitem{GC} Garc\'{\i}a Cuerva J., Weighted $H^p$ spaces,  Dissertationes Math.\textbf{ 162 }, 1-63 (1979).
		
		\bibitem{GR} Garc\'{\i}a Cuerva J., Rubio de Francia J.L., Weighted Norm Inequalities and Related Topics, North-Holland Elseviers Science publishers B.V. (1985).
		
		
		
		\bibitem{GU2} Godoy T., Urciuolo M. On certain integral operators of
		fractional type, \emph{ Acta Math. Hungar.} \textbf{82} (1-2), 99-105, (1999).
		
		
		\bibitem{G} Grafakos, L., Clasical Fourier Analysis, Second Edition, Springer Science+Business Media, LLC,  (2008).
		
		%\bibitem{GU1} Godoy T., Urciuolo M., About the $L^p$ boundedness
		%of some integral operators, \emph{Revista de la UMA}  \textbf{38}, 192--195, (1993).
		
		
		%\bibitem{JN} John, F. - Nirenberg, L., On functions of bounded mean oscilation,
		% \emph{Comm. Pure Appl. Math.} \textbf{14}, 415-426, (1961).
		
		\bibitem {KuW} Kurtz D.S., Wheeden R.L., Results on weighted norm inequalities for multipliers, \emph{Tran. Amer. Math. Soc.} \textbf{255}, 343-362, (1979).
		
		\bibitem{L} Latter, Robert. "A characterization of $ H^{p}(R^{n}) $ in terms of atoms." Studia Mathematica 62.1 (1978): 93-101.
		
		\bibitem{LeLi} Lee M., Lin C.,  The molecular characterization of weighted Hardy spaces,  \emph{J. Funct. Anal}  \textbf{188}, 442-460, (2002).
		
		\bibitem{MW} Muckenhoupt B., Wheeden R., Weighted norm inequalities for fractional integrals, \emph{Trans. Amer. Math. Soc.} \textbf{192}, 261--274, (1974).
		
		
		%\bibitem{O} R. O'Neil, {\em Fractional integration in Orlicz spaces},  Trans. Amer. Math. Soc. \textbf{115} (1963),  300--328.
		
		%\bibitem{RR} M.M. Rao and Z.D. Ren, {\em Theory of Orlicz spaces}, Monographs and Textbooks in Pure and Applied Mathematics \textbf{146}, Marcel Dekker, Inc., New York, (1991).
		
		%\bibitem{RS} Ricci, F, Sj\"{o}gren, P., Two parameter maximal functions in
		%the Heisenberg group, \emph{ Math. Z.} \textbf{199} (4), 565--575, (1988).
		
		\bibitem{RiU} Riveros, M. S.,  Urciuolo, M., Weighted inequalities for
		integral operators with some homogeneous kernels, \emph{Czech. Math. J.}, \textbf{55} (130),  423--432, (2005).
		
		\bibitem{RiU1} Riveros, M. S.,  Urciuolo, M., Weighted inequalities for fractional integral operators with some homogeneous kernels,   \emph{Acta Mathematica Sinica}, \textbf{29}(3), 449-460, (2013).
		
		\bibitem{RiU2} Riveros, M. S.,  Urciuolo, M., Weighted inequalities for some integral operators with rough kernels, \emph{Cent. Eur. J. Math}.,\textbf{  12}(4), 636-647, (2014).
		
		\bibitem{RoU}  Rocha, P.,Urciuolo, M.,  On the $H^p-L^q$ boundedness of some fractional integral operators.  \emph{Czech. Math. J}, \textbf{62} (3),  625--635, (2012).
		
		\bibitem{Rocha} Rocha P., On the atomic and molecular decomposition of weighted Hardy spaces. 
		Rev. Uni\'on Mat. Argent. \textbf{61(2)}, 229--247, (2020). 
		
		\bibitem{Rocha-2} Rocha, Pablo. "A note on Hardy spaces and bounded linear operators." Georgian Math. J 25.1 (2018): 73-76.
		
		\bibitem{Rocha-3} Rocha, Pablo. "On the atomic and molecular decomposition of weighted Hardy spaces." arXiv preprint arXiv:1710.00455 (2017).
		
		\bibitem{SW} Str\"{o}mberg, J. O., Wheeden R.L. Fractional integrals on weighted $Hp$ and $L^p$ spaces, \emph{ Trans. Amer. Math. Soc.}, \textbf{287}, 293-321, (1985).
		
		\bibitem{ST} Stromberg, Jan-Olov, and Torchinsky, A. "Weighted Hardy Spaces, vol. 1381 of Lecture Notes in Mathematics." (1989).
		
		\bibitem{Wa} Watson D.K., Weighted estimates for singular integrals via Fourier transform estimates, \emph{Duke Math. J.} \textbf{60}, 389-399, (1990).
		
		\bibitem{YZ} Yang, D. and Zhou, Y. "A boundedness criterion via atoms for linear operators in Hardy spaces." Constructive Approximation 29.2 (2009): 207-218.
		
		
		
		
	\end{thebibliography}
	
\end{document}
